\documentclass[12pt]{article}
\usepackage[margin=1in]{geometry}
\usepackage{amsmath, amssymb, amsfonts}
\usepackage{enumitem}
\usepackage{hyperref}
\setlist[enumerate]{leftmargin=*}

\title{Solutions -- Practice Problem Set 2}
\author{ECON 320 -- Introduction to Mathematical Economics}
\date{Fall 2025}

\begin{document}
\maketitle

\begin{enumerate}


\item Consider the system:
\[
\begin{cases}
x + y = \alpha \\
x - y = \beta
\end{cases}
\]

\textbf{Solution:}  
Solving:
\[
x = \tfrac{1}{2}(\alpha + \beta), \quad y = \tfrac{1}{2}(\alpha - \beta).
\]
Derivatives:
\[
\frac{\partial x}{\partial \alpha} = \tfrac{1}{2}, \quad
\frac{\partial x}{\partial \beta} = \tfrac{1}{2}, \quad
\frac{\partial y}{\partial \alpha} = \tfrac{1}{2}, \quad
\frac{\partial y}{\partial \beta} = -\tfrac{1}{2}.
\]
Interpretation: $x$ rises with both $\alpha$ and $\beta$; $y$ rises with $\alpha$ but falls with $\beta$.



\item
Let the implicit function be defined by:
\[
F(x,y,p) = x^2 + xy + y^2 - p = 0.
\]

\textbf{Solution:}  
Total differential:
\[
(2x+y)\,dx + (x+2y)\,dy - dp = 0.
\]
So:
\[
\frac{dx}{dp} = \frac{1}{2x+y}, \quad \frac{dy}{dp} = \frac{1}{x+2y}.
\]
At $(x,y)=(1,1)$:
\[
\frac{dx}{dp} = \tfrac{1}{3} > 0, \quad \frac{dy}{dp} = \tfrac{1}{3} > 0.
\]
Both variables increase with $p$.



\item
Consider the nonlinear system:
\[
\begin{cases}
x^2 + y = p \\
x + y^2 = q
\end{cases}
\]

\textbf{Solution:}  
Jacobian:
\[
J(x,y) =
\begin{bmatrix}
2x & 1 \\
1 & 2y
\end{bmatrix}, \quad
\det(J) = 4xy - 1.
\]
Condition: A locally unique solution exists if $4xy \neq 1$.



\item 
The system is:
\[
\begin{cases}
F_1(x,y,\theta) = x^2 + y - \theta = 0 \\
F_2(x,y,\theta) = xy - 1 = 0
\end{cases}
\]

\textbf{Solution:}  
Jacobian wrt $(x,y)$:
\[
J =
\begin{bmatrix}
2x & 1 \\
y & x
\end{bmatrix}.
\]
At $(1,1)$:
\[
\det(J) = 2(1) - 1 = 1 \neq 0.
\]
So IFT applies: $(x(\theta),y(\theta))$ exist and are differentiable near $\theta=2$.  

Differentiating:
\[
\begin{bmatrix}
2x & 1 \\
y & x
\end{bmatrix}
\begin{bmatrix}
dx/d\theta \\
dy/d\theta
\end{bmatrix}
=
\begin{bmatrix}
1 \\
0
\end{bmatrix}.
\]
At $(1,1)$:
\[
\begin{bmatrix}
2 & 1 \\
1 & 1
\end{bmatrix}
\begin{bmatrix}
dx/d\theta \\
dy/d\theta
\end{bmatrix}
=
\begin{bmatrix}
1 \\
0
\end{bmatrix}.
\]
Solving:
\[
dx/d\theta = 1, \quad dy/d\theta = -1.
\]



\item
Suppose:
\[
\begin{cases}
x + y + z = \theta \\
x^2 + y^2 = 1 \\
z^2 = x
\end{cases}
\]

\textbf{Solution:}  
Jacobian:
\[
J =
\begin{bmatrix}
1 & 1 & 1 \\
2x & 2y & 0 \\
-1 & 0 & 2z
\end{bmatrix}.
\]

Total differential:
\[
J
\begin{bmatrix}
dx \\ dy \\ dz
\end{bmatrix}
=
\begin{bmatrix}
d\theta \\ 0 \\ 0
\end{bmatrix}.
\]

So:
\[
\frac{\partial(x,y,z)}{\partial \theta} = J^{-1}
\begin{bmatrix}
1 \\ 0 \\ 0
\end{bmatrix}.
\]

By Cramer’s rule, these are ratios of determinants. Exact values depend on $(x,y,z)$, but the method shows how $\theta$ shifts propagate.



\item
Let:
\[
u = x^2 + y^2, \quad v = xy.
\]

\textbf{Solution:}  
Jacobian:
\[
\frac{\partial(u,v)}{\partial(x,y)} =
\begin{bmatrix}
2x & 2y \\
y & x
\end{bmatrix}, \quad
\det = 2(x^2 - y^2).
\]

If $\det \neq 0$, inverse exists:
\[
\frac{\partial(x,y)}{\partial(u,v)} =
\frac{1}{2(x^2 - y^2)}
\begin{bmatrix}
x & -2y \\
-y & 2x
\end{bmatrix}.
\]

Interpretation: Transforming to $(u,v)$ simplifies quadratic optimization by reducing cross-product terms.

\end{enumerate}
\end{document}
